\documentclass[12pt]{article}
\usepackage[utf8]{inputenc}
\usepackage[T1]{fontenc}
\usepackage[margin=1in]{geometry}
\usepackage{setspace}
\usepackage{xcolor}
\usepackage[colorlinks=true, urlcolor=blue!60!black, linkcolor=blue!60!black]{hyperref}
\setlength{\parindent}{2em}
\setlength{\parskip}{1.25em}
\renewcommand{\baselinestretch}{1.0}
\begin{document}
\begin{center}
{\LARGE \textit{Fauntleroy v. Lum}, 210 U.S. 230 (1908)\par}
\vspace{0.8em}
{\large Case Brief --- Four-Page Edition\par}
\end{center}
\vspace{1em}

\section{Memory Jogger}

A Missouri judgment on an award arising from a Mississippi cotton-futures contract that was illegal in Mississippi had to be enforced in Mississippi: the Court held 5--4, per Justice Holmes, that full faith and credit bars collateral attack on a sister-state judgment for errors of law, as distinct from a lack of jurisdiction, which remains open to collateral review (\href{https://supreme.justia.com/cases/federal/us/210/230/}{210 U.S. 230}).

\section{Detailed Case Facts}

Two Mississippi residents made a speculative cotton-futures contract in Mississippi with no intent to deliver the cotton. Under Mississippi's 1892 Code, dealing in futures was a misdemeanor, and such contracts were void and ``shall not be enforced by any court.'' When Lum refused to pay his losses, the parties submitted the dispute to arbitration in Mississippi, leaving out the question of illegality, and the arbitrator made an award against him.

\section{Procedural History}

As Justice White's dissent recounts, the claimant first sued on the award in Mississippi but dismissed once the court's attention was drawn to the contract's illegality. He then served Lum while Lum was temporarily in Missouri and sued there. The Missouri trial judge excluded Lum's evidence of illegality and told the jury to find for the plaintiff if the submission and unpaid award were proved; judgment was entered. The judgment was assigned to Fauntleroy, the claimant's attorney, who sued on it in Mississippi. The Mississippi Supreme Court upheld Lum's illegality defense, and Fauntleroy sought review by writ of error.

\section{Judicial Votes}

Reversed, \textbf{5--4} (decided May 18, 1908).
\begin{itemize}
\item \textbf{Majority:} Justice Holmes, joined by Chief Justice Fuller and Justices Brewer, Peckham, and Moody.
\item \textbf{Concurrences:} None.
\item \textbf{Dissent:} Justice White, joined by Justices Harlan, McKenna, and Day.
\end{itemize}

\section{Holding}

Mississippi's anti-futures statute set a substantive rule of decision; it did not withdraw jurisdiction from Mississippi's courts of general jurisdiction. Because the Missouri court had undisputed personal and subject-matter jurisdiction, the Full Faith and Credit Clause required Mississippi to give its judgment the same conclusive effect it had in Missouri, even if Missouri misapprehended Mississippi law and enforced an obligation illegal under Mississippi's criminal laws and public policy.

\section{Analysis of Opinions}

\textbf{Majority (Holmes): jurisdiction versus merits.} Lum relied on \textit{Anglo-American Provision Co. v. Davis Provision Co.}, where New York lawfully closed its courts to certain foreign suits. Holmes distinguished a court's power from its duty: a court of general jurisdiction may decide a case even when the law forbids recovery, and an error on the merits makes the judgment erroneous, not void. Reading ``shall not be enforced by any court'' as jurisdictional would be like treating the Statute of Frauds as jurisdictional.

\textbf{Majority (Holmes): conclusiveness.} Under \textit{Hampton v. M'Connel} and the 1790 Act, a judgment must receive the same ``credit, validity, and effect'' everywhere that it has where rendered. Contrary language in \textit{Wisconsin v. Pelican Insurance Co.} was dictum. Lum's remedy was appeal in Missouri, not collateral attack in Mississippi. Sister-state mistakes would be ``rare,'' and every State stands ``equally on both sides.''

\textbf{White (dissenting).} The ruling would ``obliterate all state lines'' by letting each State overthrow its sisters' police powers. Drawing on Story and \textit{Thompson v. Whitman}, White read the Clause as codifying comity, which never required enforcing a judgment validating a domestic crime; \textit{Pelican} allowed looking behind a judgment; the Court was bound by Mississippi's jurisdictional construction; and the arbitration was ``negligible.''

\textbf{Later cases.} \textit{Kenney v. Supreme Lodge} (1920) barred evading the rule by denying jurisdiction; \textit{Baker v. General Motors} (1998) reaffirmed that the obligation is ``exacting''; and \textit{V.L. v. E.L.} (2016) quoted the ``slow to read ambiguous words'' rule.

\section{Examples: Future Applications}

\begin{enumerate}
\item \textbf{Same-Side (Enforced): Illegal Sports Wager.} A State B judgment on a State A wager, obtained by in-state service and not appealed, must be enforced in State A despite its error about State A law.
\item \textbf{Same-Side (Enforced): Non-Compete Damages.} A State D judgment for breach of a non-compete must be enforced in State C, which voids non-competes; the remedy was appeal in State D.
\item \textbf{Opposite-Side (May Refuse): Civil Penalty.} A State B penalty against an absent defaulting defendant may be refused as penal under \textit{Huntington v. Attrill}, and personal jurisdiction may be contested.
\item \textbf{Opposite-Side (Need Not Obey): Anti-Testimony Injunction.} Under \textit{Baker}, one State's decree cannot control another's evidentiary procedures against nonparties.
\item \textbf{Fence-Sitter: Private Bounty Judgment.} A \$10,000-per-violation award to any private plaintiff looks penal yet is collected by a private party, as in \textit{Fauntleroy}.
\end{enumerate}

\section{Critique}

\textbf{Forum-shopping and asymmetry.} \textit{Fauntleroy} rewards ``tag'' jurisdiction that launders a claim barred at home into an unimpeachable sister-state judgment. Before judgment Mississippi could apply its own statute; the instant Missouri ruled, even by glaring error, that policy evaporated.

\textbf{Holmes's optimism tested.} Holmes predicted that sister-state overrides of another State's criminal or public-policy laws would be ``rare.'' Modern disputes over reproductive care, firearms liability, and non-compete bans test that assumption.

\section{Key Quotations}

\begin{quote}
``A judgment is conclusive as to all the \textit{media concludendi}, \ldots{} and it needs no authority to show that it cannot be impeached either in or out of the state by showing that it was based upon a mistake of law. Of course, a want of jurisdiction over either the person or the subject matter might be shown. \ldots{} But, as the jurisdiction of the Missouri court is not open to dispute, the judgment cannot be impeached in Mississippi even if it went upon a misapprehension of the Mississippi law.''
\end{quote}
\noindent --- Justice Holmes, 210 U.S. at 237.

\begin{quote}
``No doubt it sometimes may be difficult to decide whether certain words in a statute are directed to jurisdiction or to merits, but the distinction between the two is plain. One goes to the power, the other only to the duty, of the court.''
\end{quote}
\noindent --- Justice Holmes, 210 U.S. at 234--235.

\begin{quote}
``When it affects a court of general jurisdiction and deals with a matter upon which that court must pass, we naturally are slow to read ambiguous words as meaning to leave the judgment open to dispute or as intended to do more than to fix the rule by which the court should decide.''
\end{quote}
\noindent --- Justice Holmes, 210 U.S. at 235.

\end{document}
